\documentclass[11pt]{article}

% Core packages
\usepackage[margin=1in]{geometry}
\usepackage{amsmath,amssymb}
\usepackage{booktabs}
\usepackage{graphicx}
\usepackage{hyperref}
\usepackage{color}
\usepackage{microtype}
\usepackage{natbib}
\usepackage{times}
\usepackage{multirow}
\usepackage{array}
\usepackage{xspace}

% Annotation commands (used by /proofread-grammar and /proofread-technical)
\newcommand{\grammarcheck}[2]{\textcolor{red}{#1}\,\textcolor{red}{[#2]}}
\newcommand{\techcheck}[1]{\textcolor{blue}{[#1]}}

% Convenience macros
\newcommand{\tssmall}{T5-small\xspace}
\newcommand{\tslarge}{T5-large\xspace}
\newcommand{\gpttwo}{GPT-2\xspace}
\newcommand{\gptsmall}{GPT-2-small\xspace}
\newcommand{\gptlarge}{GPT-2-large\xspace}
\newcommand{\xaonly}{XA-only\xspace}
\newcommand{\squad}{SQuAD\xspace}
\newcommand{\hotpotqa}{HotpotQA\xspace}

\title{\textbf{Reconsidering Encoder-Decoder Architectures\\
for Context-Grounded Generation:\\
Quality, Efficiency, and Scale}}

\author{
  Mohamed El Harchaoui \\
  \texttt{med.el.harchaoui@gmail.com}
}

\date{}

\begin{document}

\maketitle

\begin{abstract}
The prevailing trend in natural language processing has shifted decisively toward
decoder-only language models. We question whether this shift is appropriate for
\emph{context-grounded generation} tasks, where the model must produce an answer
conditioned on a supplied passage rather than generating from parametric memory alone.
We conduct a controlled paradigm comparison: T5 (encoder-decoder) versus GPT-2
(decoder-only) at two parameter scales (60M/124M and 737M/774M), on two extractive
question-answering benchmarks (\squad and \hotpotqa).
T5 cross-attention-only adaptation---which trains only the decoder's cross-attention
layers (6.3M parameters, 10.4\% of the model)---outperforms fully fine-tuned GPT-2
by 2.65$\times$ on \squad and 7.0$\times$ on \hotpotqa at small scale, and by
1.61$\times$ and 3.61$\times$ at large scale.
LoRA (r=8, 0.32\% of parameters) matches full fine-tuning,
confirming that the advantage is architectural rather than a training-budget artifact.
We further show that T5's generation latency is context-length-independent
(+6\% from 64 to 512 tokens) while GPT-2's grows by 544\%, making encoder-decoder
a dominant strategy at the context lengths typical of retrieval-augmented generation.
BERTScore evaluations confirm all findings at the semantic level.
\end{abstract}

%------------------------------------------------------------------------------
\section{Introduction}
%------------------------------------------------------------------------------

Large language models based on the decoder-only architecture---GPT-2
\citep{radford2019language}, GPT-3 \citep{brown2020language}, LLaMA
\citep{touvron2023llama}---have achieved striking results on open-ended generation
tasks, and the community has widely adopted this paradigm as the default for
language modeling.
At the same time, encoder-decoder models such as T5 \citep{raffel2020exploring}
and BART \citep{lewis2020bart} have received comparatively less attention in recent
years, despite being purpose-built for sequence-to-sequence tasks.

This paper asks a precise question: \emph{for tasks that require generating a
response conditioned on a given context passage, does the encoder-decoder
architecture provide a systematic advantage over decoder-only?}
Context-grounded generation arises in retrieval-augmented generation (RAG)
\citep{lewis2020retrieval}, document-grounded dialogue, extractive QA, and
summarization---tasks that constitute a large fraction of production NLP deployments.

We argue that the answer is yes, and that the advantage is \emph{architectural}
rather than a consequence of model size, training data, or fine-tuning budget.
The key structural difference is the encoder's bidirectional attention over the
context: every context token can attend to every other context token before
the decoder generates a single answer token.
In a decoder-only model, the same context is processed with causal (left-to-right)
attention during prefill, so later context tokens attend to earlier ones but not
vice versa---a structural asymmetry that hinders tasks requiring holistic passage
comprehension.

\paragraph{Contributions.}
\begin{enumerate}
  \item We present a controlled paradigm comparison at two scales (small: 60M/124M;
        large: 737M/774M) on \squad \citep{rajpurkar2016squad} and \hotpotqa
        \citep{yang2018hotpotqa}, covering extractive and multi-hop reasoning.
  \item We show that T5 cross-attention-only adaptation (6.3M trainable parameters)
        decisively outperforms full GPT-2 fine-tuning, and that LoRA r=8
        (0.32\% of parameters) matches full T5 fine-tuning---establishing that
        architecture alignment is the binding constraint, not parameter count.
  \item We characterize the \emph{efficiency crossover}: T5 generation latency is
        context-length-independent while GPT-2's grows linearly, making
        encoder-decoder faster at context lengths beyond ${\sim}512$ tokens.
  \item We validate all results with BERTScore \citep{zhang2020bertscore},
        showing the paradigm gap is semantic, not a lexical overlap artifact.
\end{enumerate}

%------------------------------------------------------------------------------
\section{Related Work}
%------------------------------------------------------------------------------

\paragraph{Encoder-decoder models.}
T5 \citep{raffel2020exploring} unifies NLP tasks under a text-to-text framework.
BART \citep{lewis2020bart} pre-trains an encoder-decoder with denoising objectives.
Both architectures use bidirectional encoding followed by autoregressive decoding
with cross-attention to the encoder states.

\paragraph{Decoder-only models.}
GPT-2 \citep{radford2019language} and its successors use a single transformer
stack with causal self-attention.
Despite lacking an explicit encoder, decoder-only models achieve strong results on
many benchmarks by processing the context as a prefix.
GPT-3 \citep{brown2020language} demonstrated that scale can compensate for many
architectural limitations; LLaMA \citep{touvron2023llama} showed that efficient
training data curation matters more than raw parameter count.

\paragraph{Parameter-efficient fine-tuning.}
LoRA \citep{hu2022lora} adapts large models by injecting low-rank weight matrices,
typically training fewer than 1\% of parameters.
Adapter layers \citep{houlsby2019parameter} insert small bottleneck modules
between transformer layers.
Our cross-attention-only adaptation is a structurally motivated variant:
we freeze all layers except the decoder's cross-attention, which is the only
component that directly consumes encoder representations.

\paragraph{Extractive question answering.}
\squad \citep{rajpurkar2016squad} provides extractive QA over Wikipedia passages;
answers are spans within the passage.
\hotpotqa \citep{yang2018hotpotqa} requires multi-hop reasoning across two
supporting passages, testing a model's ability to compose evidence from multiple
sources.

\paragraph{Architecture comparisons.}
Several works have compared encoder-decoder and decoder-only models, but typically
at the pre-training level or on generation benchmarks where open-ended creativity
is valued \citep{tay2022ul2}.
We focus specifically on context-grounded tasks, where the input passage is the
sole source of information and the model must not hallucinate from parametric memory.

%------------------------------------------------------------------------------
\section{Models and Experimental Setup}
%------------------------------------------------------------------------------

\subsection{Architectures}

\paragraph{T5 cross-attention-only (XA-only).}
We initialize from pre-trained T5 and fine-tune \emph{only} the cross-attention
projection matrices and associated layer norms in the decoder.
All encoder weights, decoder self-attention weights, and feed-forward weights are
frozen.
This choice is structurally motivated: cross-attention is the only pathway through
which the decoder consumes encoder representations.
Training only these parameters adapts the decoder's context-reading behavior
while preserving the encoder's pre-trained bidirectional representations.
At small scale (T5-small), this yields 6.3M trainable parameters (10.4\% of 60.5M).
At large scale (T5-large), it yields 100.7M trainable parameters (13.7\% of 737M).

\paragraph{T5 LoRA (r=8).}
We apply LoRA \citep{hu2022lora} rank-8 adapters to all attention projection
matrices in the T5-large encoder and decoder.
This yields 2.4M trainable parameters (0.32\% of 737M).

\paragraph{T5 full fine-tuning.}
All 737M parameters of T5-large are updated.
We use gradient checkpointing and batch size 2 to fit in GPU memory.

\paragraph{GPT-2 full fine-tuning.}
We fine-tune all parameters of GPT-2-small (124M) and GPT-2-large (774M).
The context is prepended as a prompt: \texttt{[question] [context] [answer prefix]},
and the model is trained to predict answer tokens autoregressively.

\subsection{Datasets}

\paragraph{SQuAD.}
We use the Stanford Question Answering Dataset \citep{rajpurkar2016squad}
with 30{,}000 training examples and 512 validation examples.
The source is formatted as \texttt{question: \textit{q} context: \textit{c}};
the target is the first annotated answer span.

\paragraph{HotpotQA.}
We use \hotpotqa \citep{yang2018hotpotqa} in the same format.
HotpotQA requires connecting evidence across two passages,
making it a more demanding test of context comprehension.

\subsection{Training Details}

All models are trained with AdamW, linear warmup (10\% of steps),
bfloat16 precision, and seed 37.
We train for 3{,}000 steps with effective batch size 32 (batch 4,
gradient accumulation 8).
\techcheck{batch-size claim needs verification: T5-large used batch=2 with grad\_accum=16 (= effective 32), which differs from the stated batch=4 x accum=8; if different models used different per-GPU batch sizes, either unify the description or report per-model details to avoid misleading the reader}
Generation uses beam search with 4 beams.
Hardware: single NVIDIA RTX 3060 (12GB VRAM).

\subsection{Evaluation Metrics}

We report token-level F1 (the standard \squad metric), exact match (EM),
and BERTScore F1 \citep{zhang2020bertscore} using \texttt{roberta-large}
embeddings rescaled with the language baseline.
We also report prediction diversity (unique predictions / total) to detect
model collapse.

%------------------------------------------------------------------------------
\section{Results}
%------------------------------------------------------------------------------

\subsection{Main Quality Results}
\label{sec:main-results}

Table~\ref{tab:main-results} presents the full paradigm comparison.
At \emph{small scale}, T5-small \xaonly (60.5M total, 6.3M trainable)
outperforms fully fine-tuned \gptsmall (124M) by 2.65$\times$ on \squad
(F1: 0.708 vs.\ 0.267) and 7.0$\times$ on \hotpotqa (0.209 vs.\ 0.030).
This gap holds despite GPT-2-small having twice the total parameters.

At \emph{large scale}, T5-large \xaonly (737M total, 100.7M trainable)
outperforms \gptlarge (774M, fully fine-tuned) by 1.61$\times$ on \squad
(0.813 vs.\ 0.504) and 3.61$\times$ on \hotpotqa (0.306 vs.\ 0.085).
The parameter counts are near-matched (737M vs.\ 774M), making this a
controlled architectural comparison.

\begin{table}[t]
  \centering
  \small
  \caption{Main paradigm comparison. All models trained on 30k examples (SQuAD)
  or 30k examples (HotpotQA) except small-scale models (3k). \xaonly = cross-attention-only adaptation.
  Bold: best per task and scale group.}
  \label{tab:main-results}
  \begin{tabular}{llrrrr}
    \toprule
    \multirow{2}{*}{Model} & \multirow{2}{*}{Arch.} &
    \multirow{2}{*}{Trainable} &
    \multicolumn{2}{c}{\squad F1} & HotpotQA \\
    \cmidrule(lr){4-5}
    & & & F1 & EM & F1 \\
    \midrule
    \multicolumn{6}{l}{\emph{Small scale ($\sim$60--124M parameters)}} \\
    \gptsmall & dec-only & 124M & 0.267 & 0.150 & 0.030 \\
    \tssmall~\xaonly & enc-dec & \textbf{6.3M} & \textbf{0.708} & \textbf{0.529} & \textbf{0.209} \\
    \midrule
    \multicolumn{6}{l}{\emph{Large scale ($\sim$737--774M parameters)}} \\
    \gptlarge & dec-only & 774M & 0.504 & 0.352 & 0.085 \\
    \tslarge~\xaonly & enc-dec & 100.7M & \textbf{0.813} & \textbf{0.641} & \textbf{0.306} \\
    \tslarge~LoRA r=8 & enc-dec & 2.4M & \textbf{0.815} & 0.645 & — \\
    \tslarge~full FT & enc-dec & 737M & \textbf{0.816} & 0.660 & — \\
    \bottomrule
  \end{tabular}
\end{table}

\paragraph{Architecture dominates scale.}
A striking finding is that \tssmall (60.5M total parameters) outperforms
\gptlarge (774M parameters, 12.8$\times$ larger) on \hotpotqa
(F1: 0.209 vs.\ 0.085).
The structural capability---bidirectional encoding of multiple passages---matters
more than raw parameter count for multi-hop reasoning.

\subsection{Efficiency of Adaptation}
\label{sec:efficiency}

Table~\ref{tab:main-results} shows three T5-large training configurations
on \squad: \xaonly (13.7\% trainable), LoRA r=8 (0.32\%), and full fine-tuning
(100\%).
All three achieve F1 $\approx 0.815$, with differences smaller than measurement
noise.
Crucially, \xaonly and LoRA achieve \emph{lower} cross-entropy validation loss
than full fine-tuning (0.306 and 0.306 vs.\ 0.336), suggesting that regularization
via parameter freezing prevents mild overfitting.
\techcheck{minor precision issue: XA-only val\_loss = 0.3059 and LoRA val\_loss = 0.3062 both round to 0.306, but reporting them identically obscures that they are distinct values; consider reporting as 0.306 and 0.306 (rounded) or giving more decimal places to show they differ; also, the causal claim ``regularization prevents overfitting'' is plausible but not directly tested --- a learning curve or train-vs-val loss gap would support this interpretation}
Training only the cross-attention layers or a tiny rank-8 subspace is sufficient
to unlock the pre-trained encoder-decoder's context-reading capability.

\subsection{BERTScore: Semantic Evaluation}
\label{sec:bertscore}

Token F1 measures lexical overlap.
BERTScore \citep{zhang2020bertscore} uses contextual \texttt{roberta-large}
embeddings to capture semantic similarity.
Table~\ref{tab:bertscore} shows that the paradigm gap persists and in some cases
\emph{widens} under semantic evaluation.

\begin{table}[t]
  \centering
  \small
  \caption{BERTScore F1 (\texttt{roberta-large}, rescaled baseline).
  The paradigm gap is robust across both lexical and semantic metrics.}
  \label{tab:bertscore}
  \begin{tabular}{llrrr}
    \toprule
    Model & Task & Token F1 & BERTScore & Ratio \\
    \midrule
    \gptsmall  & \squad    & 0.267 & 0.274 & \multirow{2}{*}{2.57$\times$} \\
    \tssmall~\xaonly & \squad    & 0.707 & \textbf{0.703} & \\
    \midrule
    \gptlarge  & \squad    & 0.504 & 0.525 & \multirow{2}{*}{1.51$\times$} \\
    \tslarge~\xaonly & \squad    & 0.813 & \textbf{0.795} & \\
    \midrule
    \gptlarge  & \hotpotqa & 0.085 & 0.048 & \multirow{2}{*}{7.2$\times$} \\
    \tslarge~\xaonly & \hotpotqa & 0.306 & \textbf{0.344} & \\
    \bottomrule
  \end{tabular}
\end{table}

On \hotpotqa, GPT-2-large's BERTScore (0.048) falls \emph{below} its token F1
(0.085), while T5-large's BERTScore (0.344) \emph{exceeds} its token F1 (0.306).
This indicates that GPT-2's predictions on \hotpotqa are semantically unrelated
to the gold answers, not merely lexically imprecise---a qualitative failure mode,
not a measurement artifact.

\subsection{Long-Context Efficiency}
\label{sec:latency}

A practical concern for RAG deployments is how inference latency scales with
context length.
We benchmark single-query latency at context lengths 64--512 tokens,
holding output length constant at 16 tokens, using greedy decoding.

\paragraph{Results.}
T5-small \xaonly latency is essentially flat: 50.1~ms at 64 tokens, 53.1~ms
at 512 tokens (+6\%).
GPT-2-small latency grows from 6.3~ms to 40.6~ms over the same range (+544\%).
The curves cross at approximately 512 tokens.

\paragraph{Quality also scales better for T5.}
Evaluating both models at context lengths 128--512 tokens
(Table~\ref{tab:context-quality}), T5-small F1 increases from 0.536 to 0.722
(+35\%) while GPT-2-small increases from 0.087 to 0.288 (+20\%).
Bidirectional encoding extracts proportionally more information from longer contexts.
\techcheck{inconsistent percentage bases: ``+35\%'' for T5 is a relative increase (0.722/0.536 - 1 = 34.7\%), but ``+20\%'' for GPT-2 appears to be the absolute increase in F1 points (0.288 - 0.087 = 0.201); as a relative increase GPT-2 would be +231\%; use the same base for both, e.g.\ both as relative or both as F1-point gains}

\begin{table}[t]
  \centering
  \small
  \caption{Token F1 vs.\ context length (\squad, 512 evaluation examples).
  T5-small benefits more from longer contexts than GPT-2-small.}
  \label{tab:context-quality}
  \begin{tabular}{rrrr}
    \toprule
    Context (tokens) & \tssmall~\xaonly & \gptsmall & T5 advantage \\
    \midrule
    128 & 0.536 & 0.087 & 6.1$\times$ \\
    192 & 0.666 & 0.199 & 3.3$\times$ \\
    256 & 0.701 & 0.237 & 3.0$\times$ \\
    320 & 0.708 & 0.263 & 2.7$\times$ \\
    384 & 0.708 & 0.267 & 2.7$\times$ \\
    \textbf{512} & \textbf{0.722} & \textbf{0.288} & \textbf{2.5$\times$} \\
    \bottomrule
  \end{tabular}
\end{table}

\paragraph{Combined quality-latency tradeoff.}
At 512-token context---the typical chunk size for RAG systems---T5-small delivers
\textbf{1.94$\times$ better F1 per millisecond} and \textbf{5.2$\times$ better F1
per parameter} than GPT-2-small.

%------------------------------------------------------------------------------
\section{Analysis}
%------------------------------------------------------------------------------

\subsection{Why Encoder-Decoder Wins}

The fundamental advantage of encoder-decoder for context-grounded tasks is
\emph{bidirectional context encoding}.
In T5, every context token attends to every other context token in the encoder.
This means that when the encoder computes the representation of the answer span,
it has full information from both the question and the surrounding passage.
The decoder then reads these rich representations through cross-attention.

In GPT-2, the context is processed with causal (left-to-right) attention.
Token at position $i$ only attends to positions $j \leq i$.
This means the representation of each context token is incomplete---it cannot
attend to tokens that appear later in the passage.
For extractive QA, where the answer is a span and the question may appear before
the context, this is a structural limitation: the model cannot build a
fully-informed representation of the passage before generating the answer.

\subsection{Why the Gap Is Larger on HotpotQA}

\hotpotqa requires combining evidence from two separate passages.
The multi-hop structure demands that evidence in passage~A (about entity~X)
be combined with evidence in passage~B (about entity~Y), which requires both
passages to be mutually attended to during encoding.

In T5, both passages are concatenated and encoded bidirectionally:
the model can directly compute cross-passage attention weights, building a
joint representation of both supporting facts.
In GPT-2, the KV cache grows linearly with the concatenated input;
causal attention means later tokens (from passage~B) can attend to
earlier tokens (passage~A), but not vice versa.
This asymmetry is especially harmful for multi-hop reasoning,
explaining why the T5 advantage grows from 1.61$\times$ (SQuAD) to
3.61$\times$ (HotpotQA) at large scale.

\subsection{Why T5 Latency Is Flat}

Each T5 decoder step has cost dominated by its feed-forward sublayers
($d_\text{model} \times d_\text{ff} \approx 512 \times 2048$ multiply-accumulates
for T5-small).
\techcheck{FFN cost formula is missing the factor of 2 for the two linear projections (up- and down-projection): cost per token is $2 \times d_\text{model} \times d_\text{ff}$; also, this is cost per output token, not per step if T5 generates multiple tokens}
Cross-attention over $L$ encoder states adds $L \times d_{kv}$ work per step,
\techcheck{symbol $d_{kv}$ is never defined; the full cross-attention cost per token per step is $O(L \times d_\text{model})$ (ignoring heads); consider replacing $d_{kv}$ with $d_\text{model}$ or defining it explicitly}
which is small relative to the FFN cost for typical $L \leq 512$.
The encoder runs once (a single bidirectional pass) and is not re-invoked
during decoding.
Total latency $\approx T \times c_\text{step}$ where $T$ is the number of
generated tokens and $c_\text{step}$ is nearly context-independent.
\techcheck{$c_\text{step}$ is introduced but not defined; state explicitly that it denotes the per-step compute cost (dominated by FFN), so readers do not need to infer it}

GPT-2 with KV caching attends to all $(L + t)$ previous tokens at step $t$.
Cost per step scales as $O(L)$, making total latency $O(T \cdot L)$.
At $L = 512$, $T = 16$: this is $8{,}192$ token-attention operations
versus $\approx 512 \times 16 = 8{,}192$ cross-attention operations in T5---
but T5 computes this once (in the encoder) rather than at every decode step.
\techcheck{the numerical equality 8192 = 8192 is coincidental and potentially misleading: (1) GPT-2's total context-attention cost is $\sum_{t=0}^{T-1}(L+t) = T \cdot L + T(T-1)/2 = 8192 + 120 = 8312$, not 8192; (2) the comparison omits T5's decoder self-attention cost ($O(T^2 \cdot d)$) and GPT-2's self-attention over prompt tokens; (3) the claim that T5 ``computes this once in the encoder'' conflates encoder cross-attention with the repeated cross-attention at each decoder step --- it is the \emph{encoder self-attention} that runs once; the decoder cross-attention runs $T$ times; consider replacing this numerical argument with a cleaner asymptotic statement: T5 cross-attention cost is $O(T \cdot L)$ vs.\ GPT-2 prefix-attention cost of $O(T \cdot L)$, but T5's encoder self-attention ($O(L^2)$) runs once while GPT-2's per-step cost is $O(L+t)$}

\subsection{LoRA Matches Full Fine-Tuning}

That LoRA r=8 (2.4M parameters, 0.32\%) matches T5-large full fine-tuning
(737M parameters) on \squad (F1: 0.815 vs.\ 0.816) suggests that the
pre-trained encoder already contains most of the necessary context-reading
capability.
Fine-tuning serves primarily to redirect decoder attention toward span-extractive
outputs rather than generic language generation.
The tiny LoRA subspace is sufficient for this redirection.
This has a practical implication: an encoder-decoder model requires negligible
storage overhead per task when adapted with LoRA.

%------------------------------------------------------------------------------
\section{Conclusion}
%------------------------------------------------------------------------------

We have shown that encoder-decoder architectures provide a systematic and
substantial advantage over decoder-only architectures for context-grounded
generation, at both small and large parameter scales and across two qualitatively
different benchmarks.
The advantage is:

\begin{itemize}
  \item \textbf{Large in magnitude}: 1.61--2.65$\times$ on extractive QA,
        3.61--7.0$\times$ on multi-hop QA.
  \item \textbf{Efficient}: achievable with as few as 0.32\% of model parameters
        via LoRA, with no loss relative to full fine-tuning.
  \item \textbf{Efficiency-dominant at long context}: T5 latency is effectively
        constant over 64--512 tokens while GPT-2's grows by 544\%.
        At 512-token context, T5 delivers 1.94$\times$ better quality per
        millisecond.
  \item \textbf{Semantic}: confirmed by BERTScore (roberta-large), which reveals
        that GPT-2's failures on multi-hop QA are semantically deep, not lexical.
\end{itemize}

The decoder-only paradigm is well-suited for open-ended generation from
parametric knowledge.
For tasks that require faithfully reading and exploiting a provided context
passage---the setting that defines retrieval-augmented generation, document QA,
and grounded summarization---encoder-decoder architectures offer a more
appropriate inductive bias, greater efficiency, and stronger empirical performance.
We recommend that practitioners deploying context-grounded NLP systems
reconsider encoder-decoder models before defaulting to decoder-only.

%------------------------------------------------------------------------------
\bibliographystyle{plainnat}
\bibliography{refs}
%------------------------------------------------------------------------------

\end{document}
